% !TeX root = ../../AdvMacroNote_Sun.tex
\cleardoublepage
\pagestyle{symplain}
\thispagestyle{symplain}

\chapter*{第二版前言}
\addcontentsline{toc}{chapter}{第二版前言}
\vspace{\baselineskip}

第二版将原《高级宏观经济学 II》笔记扩展为《高级宏观经济学 I \& II》合订本。新增部分来自 2026 年秋季学期由 Yinong TAN 老师讲授的《高级宏观经济学 I》，原有内容来自 2026 年春季学期由 Chengyuan HE 老师讲授的《高级宏观经济学 II》。两部分均由 Rongqing SUN 整理。

\myskip
\mysec{课程结构}

高级宏观经济学 I 使用\ref{part:01}与\ref{part:02}，依次设置 Solow、RCK、OLG、R\&D、Romer 和人力资本模型六章。高级宏观经济学 II 使用\ref{part:03}至\ref{part:05}，保留数学工具、实际经济周期模型与新凯恩斯模型的原有内容。

\myskip
\mysec{版式说明}

两门课程共用封面、目录和全书编号，但保留各自的课程信息与视觉识别：高级宏观经济学 I 使用紫色 Accent，高级宏观经济学 II 使用橙色 Accent；页边信息和部封面会随课程自动切换。第一版前言作为项目历史记录收录在后。

\vspace{\baselineskip}
\begin{flushright}
Rongqing SUN\\
福建\hspace{0.35em}\textperiodcentered\hspace{0.35em}厦门\\
2026 年 9 月 10 日
\end{flushright}

\clearpage
\pagestyle{symplain}
